\begin{proof}
\textbf{Norm non‑negativity and degenerate cases.}  
By definition $\|v\|=\sqrt{\langle v,v\rangle}\ge 0$ for all $v\in V$.  
If $x=0$, then $\|x+y\|=\|y\|=\|0\|+\|y\|$; similarly if $y=0$. Hence the inequality holds with equality in these cases.  
Henceforth assume $x\neq0$ and $y\neq0$.

\textbf{Squaring the inequality.}  
Both $\|x+y\|$ and $\|x\|+\|y\|$ are non‑negative, so
\begin{align}
\|x+y\|\le\|x\|+\|y\|
\iff
\|x+y\|^{2}\le(\|x\|+\|y\|)^{2}.
\tag{1}
\end{align}

\textbf{Expansion of $\|x+y\|^{2}$.}  
Using linearity of the inner product in the first argument and conjugate symmetry,
\begin{align}
\|x+y\|^{2}
&=\langle x+y,\,x+y\rangle \\
&=\langle x,x\rangle+\langle x,y\rangle+\langle y,x\rangle+\langle y,y\rangle \\
&=\|x\|^{2}+2\operatorname{Re}\langle x,y\rangle+\|y\|^{2}.
\tag{2}
\end{align}

\textbf{Expansion of $(\|x\|+\|y\|)^{2}$.}
\begin{align}
(\|x\|+\|y\|)^{2}= \|x\|^{2}+2\|x\|\,\|y\|+\|y\|^{2}.
\tag{3}
\end{align}

\textbf{Cauchy–Schwarz bound.}  
The Cauchy–Schwarz inequality states
\begin{align}
|\langle x,y\rangle|\le \|x\|\,\|y\|.
\tag{4}
\end{align}
Since $\operatorname{Re}z\le|z|$ for any $z\in\mathbb{C}$,
\begin{align}
\operatorname{Re}\langle x,y\rangle\le |\langle x,y\rangle|
\le \|x\|\,\|y\|.
\tag{5}
\end{align}

\textbf{Comparison of the two sides.}  
Substituting the bound \eqref{5} into \eqref{2} yields
\begin{align}
\|x+y\|^{2}
&=\|x\|^{2}+2\operatorname{Re}\langle x,y\rangle+\|y\|^{2} \\
&\le\|x\|^{2}+2\|x\|\,\|y\|+\|y\|^{2}
=(\|x\|+\|y\|)^{2},
\end{align}
which is precisely inequality \eqref{1}.

\textbf{Taking square roots.}  
Both sides of \eqref{1} are non‑negative, so taking the (non‑negative) square root preserves the order:
\[
\|x+y\|\le\|x\|+\|y\|.
\]

\textbf{Equality condition.}  
Equality in the chain above occurs iff  

\begin{itemize}
\item $\operatorname{Re}\langle x,y\rangle = |\langle x,y\rangle|$ (so $\langle x,y\rangle$ is a non‑negative real number), and
\item equality holds in Cauchy–Schwarz, i.e.\ $x$ and $y$ are linearly dependent: there exists $\lambda\in\mathbb{C}$ with $y=\lambda x$ (or $x=0$ or $y=0$).
\end{itemize}
If $y=\lambda x$ with $x\neq0$, then $\langle x,y\rangle = \lambda\langle x,x\rangle = \lambda\|x\|^{2}$.  
The condition $\operatorname{Re}\langle x,y\rangle = |\langle x,y\rangle|$ forces $\lambda$ to be a real number and, because the real part equals the modulus, a non‑negative one. Hence $\lambda\ge0$.

Consequently,
\[
\|x+y\|=\|x\|+\|y\|
\;\Longleftrightarrow\;
\begin{cases}
x=0\ \text{or}\ y=0,\\[2pt]
\text{or } y=\lambda x \text{ with }\lambda\ge0\ (\lambda\in\mathbb{R}).
\end{cases}
\]

Thus the triangle inequality holds for all vectors in an inner‑product space, and equality occurs exactly in the stated cases. \qed
\end{proof}